\question[5]
Das folgende Programm soll das Alter abfragen und ausgeben, wie alt man nächstes Jahr ist.
Es enthält zwei Fehler, die mit Datentypen zu tun haben.

\begin{lstlisting}
alter = input('Wie alt bist du? ')
naechstes = alter + 1
print('Nächstes Jahr bist du ' + naechstes + ' Jahre alt.')
\end{lstlisting}

\begin{parts}
\part Beschreibe die beiden Fehler und gib ein korrigiertes Programm an.
\part Was passiert im korrigierten Programm, wenn der Benutzer \texttt{16,5} oder \texttt{sechzehn} eingibt?
\end{parts}

\begin{solutionbox}{7cm}
\begin{parts}
\part \texttt{input()} liefert immer einen String. In Zeile 2 kann man zu einem String keine Zahl addieren
(\texttt{TypeError}), der String muss zuerst mit \texttt{int()} umgewandelt werden.
In Zeile 3 wird ein String mit einer Zahl verkettet (\texttt{TypeError}). Python wandelt die Zahl nicht automatisch
in einen String um, man braucht \texttt{str()} oder einen f-String.
\begin{lstlisting}
alter = int(input('Wie alt bist du? '))
naechstes = alter + 1
print('Nächstes Jahr bist du ' + str(naechstes) + ' Jahre alt.')
\end{lstlisting}
\part In beiden Fällen kann \texttt{int()} den String nicht in eine ganze Zahl umwandeln. Das Programm bricht mit
einem \texttt{ValueError} ab.
\end{parts}
\end{solutionbox}
